\bookchapter{פרק 3：חיבור הטלפון ל“קו החיים”——האינטרנט}

\chapterintro{ללא אינטרנט, הטלפון החכם “מתפגר” בחצי. האינטרנט הוא קו החיים של הטלפון המודרני, להבין אותו, להשתמש בו נכון ולהגן עליו – זהו שיעור חובה לכל משתמש טלפון.}

\booksection{אל תתבלבלו בחיבור——הבחנה בין “הדרך” ל“רכב” של Wi-Fi וסלולר}

דמיינו את האינטרנט כעיר.\textbf{“הדרך”}היא הנתיב מהטלפון שלכם לאינטרנט,

\begin{tcbinfo}{"הרכב"}
הם הנתונים שנוסעים על הדרך.Wi-Fi וסלולר הן שתי דרכים שונות, אך הנתונים (הרכב) עצמם זהים.\textbf{1. ההבדל המהותי:ערוץ מקומי מול כביש מהיר מרוחק}Wi-Fi היא בעצם\textbf{טכנולוגיית רשת מקומית}, כמו כביש סלול אלחוטי מתחת לבניין שלכם, האחראי על העברת נתונים מהנתב לטלפון. אבל אם הפס הרחב מנותק או הסיב האופטי נקרע, הכביש עדיין שם אבל “הכביש המהיר” החיצוני מנותק, ומטבע הדברים אי אפשר להגיע לשום מקום. לעומת זאת, תעבורת הרשת (נתונים סלולריים) עוברת דרך תחנות בסיס 4G/5G של הספק, כלומר אתם נוסעים ישירות על הנתיב המהיר של הספק, וכל עוד יש קליטה אפשר להתחבר ישירות לאינטרנט.

\end{tcbinfo}

\booksubsection{עלויות ותרחישים:תשלום חודשי קבוע מול חיוב לפי נפח}

Wi-Fi היא בדרך כלל השקעה חד-פעמית (דמי מנוי חודשיים לפס רחב), כל עוד החשמל לא נפסק, כל המשפחה יכולה להשתמש ללא הגבלה, מתאים מאוד להורדת סרטים, משחקים ותרחישי תעבורה כבדה. לעומת זאת, נתונים סלולריים מחויבים לפי נפח שימוש, חריגה מהחבילה עלולה לגרום להאטה או חיוב נוסף, מתאים יותר למשימות קלות בתרחישי ניידות כמו רכבת תחתית או נסיעות עסקים.\textbf{בואו ננסה:בדקו את חבילת הגלישה שלכם}פתחו את [הגדרות] > [רשת סלולרית] > [ניהול תעבורה] בטלפון, ובדקו כמה נפח גלישה נשאר לכם החודש ואילו אפליקציות צורכות הכי הרבה. כשיודעים מה המצב, יוצאים החוצה בראש שקט.

\booksubsection{שיקולי אבטחה:Wi-Fi ציבורי הוא “שדה דיג” נסתר}

מבחינת אבטחה, נתונים סלולריים בדרך כלל בטוחים יותר מ-Wi-Fi ציבורי. נתונים סלולריים מוצפנים על ידי הספק, בעוד ש-Wi-Fi ציבורי אכן חשוף לכך שהאקרים יקימו “נקודה חמה מזויפת” ויירטו תעבורה לא מוצפנת. אבל צריך להבחין כאן:בעידן של היום בו HTTPS נפוץ בכל מקום, מה שהתוקף יכול ליירט זה

\begin{tcbpractice}{באילו אתרים ביקרתם}
(שמות דומיין), וכן

\end{tcbpractice}

\booksubsection{תוכן HTTP לא מוצפן}

;ואילו הודעות WeChat, סיסמאות כניסה לאפליקציות בנק ונתונים המוצפנים ב-HTTPS, התוקף\textbf{לא יכול לקרוא ישירות}——אלא אם יבצע התקפת SSL-stripping מורכבת יותר והמשתמש יתעלם מאזהרת התעודה. אולם גם כך, בעת ביצוע פעולות תשלום בחוץ, שימוש ישיר בנתונים סלולריים הוא עדיין הבחירה הבטוחה ביותר.\textbf{תזכורת אבטחה}בשדות תעופה, קניונים, בתי קפה ומקומות ציבוריים אחרים, אל תתחברו ל-Wi-Fi חינמי עם שם חשוד (כמו “Free\_WiFi\_2024”). גם אם התחברתם, לעולם אל תיכנסו לאפליקציית בנק או תזינו סיסמה דרכו.\textbf{4. פענוח בטלפון:הבינו את “שפת הסתרים” של שורת המצב}איך יודעים באיזו רשת הטלפון משתמש כרגע? בטלפוני Huawei, למשל, הופעת

\textbf{אם הגשת בטעות מידע רגיש, לא מספיק למחוק את הקובץ (הוא עדיין בהיסטוריה), צריך להשתמש ב-}\par

ליד אות הספק או ליד אות ה-WLAN בשורת המצב, משמעה שהטלפון משתמש בשיטה המתאימה לגלישה;אם שניהם דולקים בו-זמנית, פירוש הדבר שהאצת רשת פועלת, והטלפון מנצל את שני הערוצים בו-זמנית.

\booksubsection{אצל משתמשי iPhone זה פשוט יותר:אם סמל ה-Wi-Fi דולק, הטלפון משתמש ב-Wi-Fi, כשלא – הביטו בפסי האות בפינה הימנית העליונה, זה סימן שאתם משתמשים בנתונים סלולריים. בחלק מטלפוני האנדרואיד יופיע גם “שירותי מיקום” או “4G” בשורת המצב, המציין ישירות את סוג הרשת הסלולרית הנוכחית.}

א. רשת סלולרית (4G/5G) ותעבורת נתונים\textbf{רשת סלולרית היא הבסיס לגלישה בכל מקום ובכל זמן. מ-2G ועד 5G, המהירות גדלה פי אלפי מונים, אבל העקרון דומה.}תקני רשת:

4G LTE

\booksection{א. רשת סלולרית (4G/5G) ותעבורת נתונים}

5G

\booksubsection{:מהיר (תיאורטית עד 10Gbps), אבל הכיסוי עדיין מתרחב, צריכת חשמל גבוהה. מתאים להורדת קבצים גדולים וצפייה בווידאו HD.}

\begin{itemize}
  \item \textbf{5G SA מול NSA}:SA היא רשת עצמאית, NSA מסתמכת על תחנות 4G. SA מספקת חוויית משתמש טובה יותר, אך דורשת תמיכה גם מהספק וגם מהטלפון.
  \item \textbf{שירותי מיקום}מהי תעבורת נתונים
  \item \textbf{:זוהי כמות הנתונים שהטלפון מייצר בגלישה. צפייה בווידאו וגלילת תמונות צורכות הכי הרבה, טקסט וקול חוסכים יחסית.}טיפים לחיסכון בתעבורה
\end{itemize}

\booksubsection{:הפעילו “מצב חיסכון בנתונים” בהגדרות, עדכנו אפליקציות וצפו בווידאו HD רק ב-Wi-Fi.}

\begin{itemize}
  \item \textbf{פסי קליטה $\neq$ מהירות}:קליטה חזקה רק מעידה על חיבור טוב, המהירות תלויה גם בעומס על תחנת הבסיס ובמרחק שלכם ממנה.
  \item \textbf{אזהרה חשובה:כבו את [נדידת נתונים] לפני נסיעה לחו"ל!}שימוש בנתונים סלולריים בחו"ל עלול לייצר חשבון עתק! לפני נסיעה לחו"ל יש לכבות את נדידת הנתונים בהגדרות, או לרכוש מראש כרטיס SIM מקומי/חבילת נדידה בינלאומית.
  \item \textbf{ב. Wi-Fi:חיבור, שיתוף ואבטחה}Wi-Fi הוא כוח הגלישה העיקרי בבית ובמשרד, מהיר ולא צורך תעבורה סלולרית. אבל גם ה-Wi-Fi מתפתח כל הזמן, עם הבדלים גדולים בין תקנים שונים.
\end{itemize}

\textbf{התחברות ל-Wi-Fi:}\par

גשו ל-[הגדרות] > [WLAN] בחרו רשת והזינו סיסמה. כיום טלפונים רבים תומכים בסריקת קוד QR להתחברות, נוח יותר.

\booksection{ב. Wi-Fi：חיבור, שיתוף ואבטחה}

גם אנדרואיד וגם iOS מאפשרים לטלפון מחובר לשתף במהירות את הסיסמה עם חברים. אנדרואיד מייצר קוד QR, iOS יכול לשתף אוטומטית בקרבה.

\booksubsection{התפתחות תקני Wi-Fi:}

תקני Wi-Fi נקבעים על ידי IEEE, כל דור מביא שיפור משמעותי:

\begin{center}
\begin{tabularx}{\linewidth}{|X|X|X|X|X|}
\hline
\textbf{תקן} & \textbf{פרוטוקול} & \textbf{תדר} & \textbf{מהירות מקסימלית} & \textbf{מאפיינים} \tabularnewline
\hline
Wi-Fi 4 & 802.11n & 2.4/5GHz & 600Mbps & תקן ישן, כיסוי סביר \tabularnewline
Wi-Fi 5 & 802.11ac & 5GHz & 3.5Gbps & הנפוץ ביותר כיום, מהירות מספקת \tabularnewline
Wi-Fi 6 & 802.11ax & 2.4/5GHz & 9.6Gbps & ביצועים מעולים בריבוי מכשירים, השהיה נמוכה \tabularnewline
Wi-Fi 6E & 802.11ax & 2.4/5/6GHz & 9.6Gbps & מוסיף תדר 6GHz, ללא עומס \tabularnewline
Wi-Fi 7 & 802.11be & 2.4/5/6GHz & 46Gbps & העברה במהירות על, צירוף קישורים מרובים \tabularnewline
\hline
\end{tabularx}
\end{center}

\begin{itemize}
  \item \textbf{טכנולוגיות מפתח של Wi-Fi 6}:OFDMA (חלוקת תדרים אורתוגונלית) מאפשרת למספר מכשירים לשדר בו-זמנית, MU-MIMO (מספר משתמשים, קלט ופלט מרובים) משפרת את יעילות העבודה הסימולטנית של מספר מכשירים, TWT (זמן התעוררות ממוקד) חוסך חשמל במכשיר.
  \item \textbf{Wi-Fi 6E}:מוסיף תדר 6GHz על בסיס Wi-Fi 6, מספק יותר ערוצים ומפחית הפרעות. דורש תמיכה גם בטלפון וגם בנתב.
  \item \textbf{Wi-Fi 7}:מציג MLO (Multi-Link Operation), מאפשר שימוש במספר תדרים בו זמנית להעברת נתונים, מפחית השהיה משמעותית. ערוץ רוחב 320MHz מביא לקפיצת מהירות. פלטפורמות שבבי דגל כמו Snapdragon 8 Elite Gen5 ו-Dimensity 9400 כבר תומכים ב-Wi-Fi 7. Xiaomi 17, Huawei Mate 80, Samsung S26 ומכשירים נוספים יכולים להשתמש ברשת Wi-Fi 7 (נדרש נתב Wi-Fi 7).
  \item \textbf{2.4GHz מול 5GHz}:2.4GHz חודר קירות טוב אבל איטי ורווי הפרעות; 5GHz מהיר אבל חודר קירות חלש. תדר 6GHz של Wi-Fi 6E/7 הוא המהיר ביותר, אך חודר קירות הכי חלש, מתאים לשימוש באותו חדר.
\end{itemize}

\booksubsection{סיכוני Wi-Fi ציבורי:}

:Wi-Fi ציבורי עלול לאפשר להאקרים ליירט נתונים או להציב נקודות חמות מזויפות. הכללים:

\begin{itemize}
  \item הימנעו מכניסה לחשבון בנק או הזנת סיסמאות.
  \item עדיף להשתמש בחיבור VPN מוצפן.
  \item אל תתחברו אוטומטית ל-Wi-Fi לא מוכר.
  \item השתמשו באתרי HTTPS (סמל מנעול בשורת הכתובת).
\end{itemize}

\begin{itemize}
  \item ג. Bluetooth:השחקן הרב-תכליתי של תקשורת אלחוטית לטווח קצר
  \item Bluetooth היא כנראה אחת הטכנולוגיות האלחוטיות הנפוצות ביותר בטלפון שלכם——אוזניות אלחוטיות, שעונים חכמים, מערכות רכב, מקלדות ועכברים, כולם מתחברים דרכה.
  \item עקרונות טכניים:
  \item Bluetooth פועלת בתדר ISM של 2.4GHz, ומשתמשת בטכנולוגיית דילוג תדרים (FHSS) לפיזור נתונים על פני 79 ערוצים במעבר מהיר, מה שמונע הפרעות ביעילות. טווח התקשורת של Bluetooth מתחלק לשלוש רמות:Class 1 כ-100 מטר, Class 2 כ-10 מטר (הנפוץ בטלפונים), Class 3 כ-1 מטר.
\end{itemize}

\booksection{ג. Bluetooth：השחקן הרב-תכליתי של תקשורת אלחוטית לטווח קצר}

Bluetooth 4.0

\booksubsection{תרחישי יישום:}

Bluetooth 4.2

\booksubsection{:שיפור מהירות העברה והגנת פרטיות, תמיכה בחיבור ישיר IPv6.}

\begin{itemize}
  \item \textbf{Bluetooth 5.0}:מהירות העברה כפולה (2Mbps), טווח פי 4 (תיאורטית 300 מטר), יכולת שידור פי 8. תמיכה בשמע כפול (חיבור לשתי אוזניות בו-זמנית).
  \item \textbf{Bluetooth 5.1}:הצגת פונקציית איתור כיווני (AoA/AoD), דיוק מיקום ברמת סנטימטרים.
  \item \textbf{Bluetooth 5.2}:תקן LE Audio חדש, תמיכה בקידוד LC3 (איכות שמע טובה יותר, צריכת חשמל נמוכה יותר), שמע רב-זרמי, שמע משודר.
  \item \textbf{Bluetooth 5.3/5.4}:אופטימיזציה נוספת של צריכת חשמל ויציבות חיבור, שידור מחזורי משופר.
  \item \textbf{מצב האימוץ של LE Audio}:בשנים 2024-2025, LE Audio נכנס לשלב אימוץ מהיר. אוזניות True Wireless חדשות של מותגים מובילים (כמו Huawei FreeBuds Pro 4, Xiaomi Buds 5 Pro, Samsung Galaxy Buds3 וכו') תומכות כולן ב-LE Audio וקידוד LC3, ומספקות השהיה נמוכה יותר ואיכות שמע טובה יותר. בקניית אוזניות חדשות, מומלץ להעדיף מוצרים התומכים ב-LE Audio.
  \item \textbf{תרחישי שימוש:}העברת שמע
\end{itemize}

\textbf{:אוזניות אלחוטיות, מערכות שמע לרכב. LE Audio של Bluetooth 5.2 מביא קפיצת מדרגה באיכות השמע וההשהיה.}:ב-2025-2026, LE Audio נכנס לשלב הפצה מהירה. אוזניות אלחוטיות חדשות של מותגים מובילים (כמו Huawei FreeBuds Pro 5, Xiaomi Buds 6 Pro, Samsung Galaxy Buds3 Pro) כבר תומכות ב-LE Audio ובקידוד LC3, עם השהיה נמוכה יותר ואיכות צליל טובה יותר. בקניית אוזניות חדשות, מומלץ לבחור מוצר התומך ב-LE Audio.

\booksubsection{: השהיה נמוכה יותר, סנכרון אודיו-וידאו במשחקים.}

\begin{itemize}
  \item \textbf{צימוד מכשירים}:שעונים חכמים, צמידים, מקלדות, עכברים, בקרי משחק.
  \item \textbf{שירותי מיקום}:Beacons של Bluetooth לניווט פנימי ומעקב אחר חפצים (כמו Apple AirTag).
  \item \textbf{בית חכם}:רשת Bluetooth Mesh לשליטה בנורות, מנעולים וכו'.
  \item \textbf{ד. שלט אינפרא-אדום:הטלפון הופך לשלט אוניברסלי}שלט אינפרא-אדום היא אחת מטכנולוגיות השליטה האלחוטיות הוותיקות ביותר, אך עדיין שימושית בטלפונים. חלק מטלפוני האנדרואיד (כמו דגמים מסוימים של Xiaomi ו-Huawei) כוללים משדר אינפרא-אדום מובנה, ויכולים לשלוט בטלוויזיות, מזגנים, ממירים ומכשירי חשמל ביתיים אחרים.
  \item \textbf{קישוריות לרכב}שלט אינפרא-אדום משתמש בדיודה פולטת אור אינפרא-אדום כדי לפלוט פולסי אור בתדר מסוים (בדרך כלל 38kHz), ומקלט האינפרא-אדום של מכשיר החשמל מפענח פולסים אלה ומבצע את הפעולה המתאימה. מותגים שונים משתמשים בפרוטוקולי קידוד שונים, הנפוצים הם NEC, RC5, RC6, SONY וכו'. אפליקציית האינפרא-אדום בטלפון מכילה ספריית קידודים של מותגים שונים, בוחרים מותג ודגם ומתאימים.
\end{itemize}

\booksection{ד. שלט אינפרא-אדום：הטלפון הופך לשלט אוניברסלי}

כיווניות

\booksubsection{דגמים נתמכים:}

מרחק קצר

\booksubsection{:המרחק האפקטיבי הוא בדרך כלל 5-8 מטר.}

\begin{itemize}
  \item \textbf{אפל לא תומכת}:ל-iPhone אין משדר אינפרא-אדום, לא ניתן להשתמש בפונקציה זו.
  \item \textbf{5. NFC:הקסם של נגיעה-ונגיעה}NFC(Near Field Communication,תקשורת טווח קצר)היא טכנולוגיית תקשורת אלחוטית בתדר גבוה לטווח קצר במיוחד(עד 10 ס"מ),תדר עבודה 13.56MHz,קצב העברה מרבי 424Kbps.
  \item \textbf{עקרונות תקשורת ומצבים:}NFC מבוססת על טכנולוגיית RFID,תומכת בשלושה מצבי עבודה:
\end{itemize}

\booksection{NFC：הקסם של נגיעה-ונגיעה}

:הטלפון קורא מידע מתגי NFC,כגון קריאת מדבקות NFC להפעלת פעולות.

\booksubsection{מצב עמית-לעמית}

:שני מכשירי NFC מעבירים נתונים ביניהם,כגון Android Beam(הופסק).

\begin{itemize}
  \item \textbf{מצב הדמיית כרטיס}:הטלפון מדמה כרטיס חכם,זהו המצב הנפוץ ביותר——כרטיס תחבורה ציבורית,כרטיס כניסה,כרטיס בנק כולם במצב זה.
  \item \textbf{מנגנוני אבטחה:}טווח קצר במיוחד
  \item \textbf{:מרחק התקשורת של עד 10 ס"מ הוא כשלעצמו אבטחה פיזית,ציתות מרחוק כמעט בלתי אפשרי.}שבב אבטחה
\end{itemize}

\booksubsection{:הטלפון מכיל שבב אבטחה עצמאי(SE)או משתמש ב-HCE(הדמיית כרטיס מארח),מפתחות התשלום מאוחסנים באזור מאובטח,גם אם הטלפון עבר root לא ניתן לחלץ אותם.}

\begin{itemize}
  \item \textbf{אימות טביעת אצבע/זיהוי פנים}:התשלום דורש אימות ביומטרי,גם אם הטלפון אבד לא ניתן לגנוב כסף.
  \item \textbf{טוקניזציה}:Apple Pay, Huawei Pay וכו' לא מעבירים מספר כרטיס אמיתי,אלא משתמשים באסימון חד-פעמי,גם אם יורט לא ניתן לשימוש חוזר.
  \item \textbf{יישומים אופייניים:}כרטיס תחבורה ציבורית/כרטיס כניסה
  \item \textbf{:כתוב את פרטי הכרטיס הפיזי לטלפון,הצמד את הטלפון וזהו.}תשלום נייד
\end{itemize}

\booksubsection{:Apple Pay, Huawei Pay וכו',קשר כרטיס בנק לטלפון,הצמד לשלם.}

\begin{itemize}
  \item \textbf{התאמת מכשירים}:NFC הצמד-הצמד להתאמה מהירה של אוזניות בלוטות',רמקולים.
  \item \textbf{אימות זהות}:קריאת תעודת זהות ב-NFC,אימות רישיון נהיגה אלקטרוני.
  \item \textbf{בית חכם}כשיש לך הרבה נפח גלישה,תוכל להפעיל נקודה חמה לטאבלט,מחשב.
  \item \textbf{איך להפעיל}:הגדרות > נקודה חמה אישית/שיתוף רשת,הפעל.
\end{itemize}

\booksection{נקודה חמה אישית：הפוך את הטלפון לנתב Wi-Fi}

:צורך הרבה סוללה ונפח גלישה,זכור להגדיר סיסמה מורכבת,ולהגביל את מספר המכשירים המחוברים.

\begin{itemize}
  \item \textbf{שיתוף USB}:חבר למחשב בכבל נתונים,יציב ומהיר יותר מנקודה חמה אלחוטית.
  \item \textbf{הערות חשובות}פתח [הגדרות],מצא "נקודה חמה אישית" או "שיתוף רשת",הפעל,הגדר שם וסיסמה.
  \item \textbf{חפש במכשיר אחר(למשל מחשב)את ה-Wi-Fi הזה והתחבר.}בדוק בצד הטלפון,האם אתה יכול לנהל מכשירים מחוברים,או לנתק מכשירים לא מוכרים.
\end{itemize}

\begin{tcbpractice}{פרק 4:התקנה וניהול אפליקציות——"חבילת הכישורים" של הטלפון}
\begin{enumerate}
  \item אפליקציות הופכות את הטלפון לכל-יכול,אבל צריך גם לנהל אותן היטב.בדיוק כמו שהיית מקצה מקום עבודה והרשאות לעובד חדש,התקנת אפליקציות דורשת גם תהליך מסודר.
  \item 1. חנות רשמית:"סופרמרקט האפליקציות" הבטוח ביותר
  \item Android
\end{enumerate}

\end{tcbpractice}
